\section{Calibrated base closure and passage beyond the Poisson length}
\label{sec:tpc27-passage-theorem}

We can now state the complete selected-base theorem and splice it to the
annular cutoff result of TPC-26.  Every conclusion concerns the same
physical row family, mask, residue factor, and smooth orbit weight at
the two cutoffs.

\subsection{The calibrated base theorem}

\begin{theorem}[Physical calibrated base closure]
\label{thm:tpc27-calibrated-base}
Retain all hypotheses and definitions of
\cref{sec:tpc27-interface-bridge,sec:tpc27-multiple-calibration}, in
particular \(R\le S\), the actual opened-row coefficients, and the
divisor-independent off-diagonal mask.  Assume also the inherited
TPC-20--23 primitive, no-wrap, Fourier-tail, and fixed-margin
interfaces.  Write the scales as in \eqref{eq:tpc27-base-exponents}.
If
\begin{equation}
 s<1-\beta,
 \qquad
 \eta_S(\beta,s)>0
 \label{eq:tpc27-base-hypotheses}
\end{equation}
with fixed margins, then for every fixed \(B>0\),
\begin{equation}
 \boxed{
 \cQ_S^{\rm phys}
 \ll_{B,\mathscr D}XQ(\log X)^{-B}.}
 \label{eq:tpc27-calibrated-base}
\end{equation}
Here \(\cQ_S^{\rm phys}\) is the calibrated physical packet
\eqref{eq:tpc27-physical-base}; the theorem is not merely an estimate
for its separated nonzero conductor cells.
\end{theorem}

\begin{proof}
Use the exact decomposition
\eqref{eq:tpc27-ZZ-EZ-ZE-EE}.  For \(ZZ\), first extract physical
additive frequency zero before smooth separation and apply
\cref{thm:tpc27-physical-base-zero}.  Then separate only the nonzero
spectrum and apply \cref{thm:tpc27-physical-nonzero}.  The single and
constant nonzero channels removed by the TPC-20 two-scale completion,
the dynamic Ramanujan mean, and the physical Fourier tail are already
included in that theorem's reassembly.  Finally apply
\cref{thm:tpc27-scalar-closure} to \(EZ,ZE,EE\).  All fixed dyadic,
sign, content, residue, and smooth decompositions cost at most a fixed
power of \(\log X\), which is absorbed by choosing the calibration and
tail orders after \(B\).
\end{proof}

\begin{remark}[Static base]
\label{rem:tpc27-static-base}
At \(S=R=X^{1/2-\delta+o(1)}\), the proof uses the TPC-22 static
theorem for the nonzero component.  Its conductor saving is
\[
 \eta_*(\beta,\delta)=
 \min\left\{
 1-\frac{3\beta}{2},
 \frac{\beta+4\delta-1}{2},
 \frac{6\delta-\beta}{4}
 \right\},
\]
which is exactly \(\eta_S(\beta,s)\) after
\(s=1/2-\delta\).  The additive-zero and scalar bridges in the present
paper are still required; the static conductor theorem alone does not
display them.
\end{remark}

\subsection{Splicing to an upper cutoff}

Let \(T=X^{t+o(1)}\) with \(S<T\).  TPC-26 proves, under its annular
support and symmetric exact-conductor hypotheses, that
\begin{equation}
 \cQ_T^{\rm phys}-\cQ_S^{\rm phys}
 \ll_{B,\mathscr D}XQ(\log X)^{-B}
 \label{eq:tpc27-imported-annular-stability}
\end{equation}
whenever its sharp relaxed-box minimax \(M_\beta(t)\) is positive with
a fixed margin \citep{WangTPC26}.  For \(1/2\le\beta<1\),
\begin{equation}
 M_\beta(t)=
 \begin{cases}
 1-\beta-t,
 &0\le t\le(1-\beta)/2,\\[2pt]
 (3-3\beta-2t)/4,
 &(1-\beta)/2\le t\le\beta/2,\\[2pt]
 (3-\beta-6t)/4,
 &\beta/2\le t\le(3\beta-1)/2,\\[2pt]
 (\beta+1-4t)/2,
 &t\ge(3\beta-1)/2.
 \end{cases}
 \label{eq:tpc27-Mbeta}
\end{equation}

\begin{corollary}[Selected calibrated passage]
\label{cor:tpc27-selected-passage}
Assume the hypotheses of \cref{thm:tpc27-calibrated-base}, with
\(1/2\le\beta<1\).  Suppose also that \(R\le S<T\) lies in the
TPC-26 annular interface and
\begin{equation}
 M_\beta(t)>0
 \label{eq:tpc27-annular-positive}
\end{equation}
with a fixed margin.  Then, for every fixed \(B>0\),
\begin{equation}
 \boxed{
 \cQ_T^{\rm phys}
 \ll_{B,\mathscr D}XQ(\log X)^{-B}.}
 \label{eq:tpc27-upper-closure}
\end{equation}
The upper cutoff \(T\) may exceed the Poisson length \(J\).
If, in addition, \(s<1-\beta<t\), this is a genuine passage from a
Poisson-short base to a Poisson-long upper cutoff.
The row family, both coefficient sequences, mask, periodic residue
factor, smooth orbit weight, and prefix definition are the same at
the two cutoffs.
\end{corollary}

\begin{proof}
Use the exact identity
\[
 \cQ_T^{\rm phys}
 =\cQ_S^{\rm phys}
 +(\cQ_T^{\rm phys}-\cQ_S^{\rm phys})
\]
and combine \eqref{eq:tpc27-calibrated-base} with
\eqref{eq:tpc27-imported-annular-stability}.
\end{proof}

The conclusion is stronger than cutoff stability alone: the base value
is now estimated in the overlap with the inherited Poisson-short
region.  It remains a selected truncated-square theorem and does not
restore the complement beyond \(T\).

\subsection{An exact rational passage sample}

Take
\begin{equation}
 \delta=\frac3{20},
 \qquad
 \beta=\frac{31}{50},
 \qquad
 s=\frac7{20},
 \qquad
 t=\frac{39}{100}.
 \label{eq:tpc27-strict-sample}
\end{equation}
Then
\begin{equation}
 R=S=X^{7/20+o(1)},
 \qquad
 J=X^{19/50+o(1)},
 \qquad
 T=X^{39/100+o(1)},
 \label{eq:tpc27-strict-order}
\end{equation}
so
\begin{equation}
 S<J<T.
 \label{eq:tpc27-crossing-order}
\end{equation}
The strict Poisson-short margin is
\begin{equation}
 (1-\beta)-s
 =\frac3{100}.
 \label{eq:tpc27-single-margin}
\end{equation}
The three base conductor entries are
\begin{equation}
 \left(
 1-\frac{3\beta}{2},
 \frac{\beta+1-4s}{2},
 \frac{3-\beta-6s}{4}
 \right)
 =\left(\frac7{100},\frac{11}{100},\frac7{100}\right),
 \label{eq:tpc27-base-sample-ledger}
\end{equation}
so the centered nonzero conductor ledger has margin \(7/100\).  The complete
calibrated base conclusion remains the arbitrary-logarithmic estimate
\eqref{eq:tpc27-calibrated-base}.  Since
\begin{equation}
 \frac\beta2=\frac{31}{100}
 <\frac{39}{100}
 <\frac{3\beta-1}{2}=\frac{43}{100},
\end{equation}
the third line of \eqref{eq:tpc27-Mbeta} applies and gives
\begin{equation}
 M_{31/50}(39/100)
 =\frac{3-31/50-6(39/100)}4
 =\frac1{100}.
 \label{eq:tpc27-annular-sample-saving}
\end{equation}
Thus the annular minimax has margin \(1/100\); as usual, the estimates
may use any fixed exponent strictly below the displayed ledger margin.

\begin{corollary}[Closed strict crossing sample]
\label{cor:tpc27-strict-sample}
For the actual opened-row packet with the fixed interfaces above, the
parameters \eqref{eq:tpc27-strict-sample} satisfy
\begin{equation}
 \cQ_T^{\rm phys}
 \ll_{B,\mathscr D}XQ(\log X)^{-B}
 \label{eq:tpc27-strict-sample-closure}
\end{equation}
for every fixed \(B>0\).  This upgrades the corresponding TPC-26
sample from a conditional base splice to a closed selected calibrated
truncated square.
\end{corollary}

\subsection{The high-\texorpdfstring{\(\beta\)}{beta} boundary}

TPC-26 also records
\begin{equation}
 \beta=\frac{267}{400},
 \qquad
 s=\frac{23}{60},
 \qquad
 t=\frac{193}{500}.
 \label{eq:tpc27-high-sample}
\end{equation}
Here
\begin{equation}
 s-(1-\beta)
 =\frac{23}{60}-\frac{133}{400}
 =\frac{61}{1200}>0.
 \label{eq:tpc27-high-base-beyond-J}
\end{equation}
Thus the base already lies beyond \(J\), and
\eqref{eq:tpc27-base-hypotheses} fails.  The annular theorem still
proves cutoff stability between \(S\) and \(T\), but the present paper
does not supply its initial value.  No high-\(\beta\) base closure is
claimed.

\subsection{What remains}

After \cref{cor:tpc27-selected-passage}, three independent walls remain:
\begin{enumerate}[label=\textup{(\roman*)}]
 \item a base nonzero theorem when the lower cutoff itself is beyond
 the Poisson length, needed for samples such as
 \eqref{eq:tpc27-high-sample};
 \item the ultra-long complement beyond the selected upper cutoff,
 whose divisor reflection produces polynomial families of two-point
 affine M\"obius correlations;
 \item the separate positivity and main-term problem required to pass
 from an analytic upper/cancellation estimate to a prime-pair
 asymptotic or lower bound.
\end{enumerate}
The first two are analytic cancellation gates.  The third remains
parity-sensitive.  None is resolved by the calibrated base theorem.
